% GENERATED FILE. Edit canonical lessons, then run npm run generate:books.
% canonical-source-hash: fnv1a64:3f1583b1ac4d43f5
% canonical-lessons: TE-C508-pogaru, TE-C508-monditanam, TE-C508-pattudala, TE-C508-sahasam, TE-C508-piriki

\chapter{Stubborn or Brave}
\label{ch:te-stubborn-or-brave}
\hlchaptermodality{508}

\begin{chapteropening}
\textbf{By the end of this chapter:} \emph{I can say that someone is cheeky, stubborn, determined, bold or timid.}

Complete the last lesson of chapter 508: I can say that someone is cheeky, stubborn, determined, bold or timid.
\end{chapteropening}

\section[\texorpdfstring{\te{పొగరు} (pogaru)}{పొగరు (pogaru)}]{\te{పొగరు} (pogaru) — conceit; cheek}

\label{lesson:TE-C508-pogaru}



\begin{quote}
Before the new one: say the Telugu for humility, then the Telugu for arrogance.
\end{quote}

\subsection*{You'll want to know: \te{పొగరు}}
\textbf{\te{పొగరు}} — \emph{pogaru} — \textquotedblleft{}conceit; cheek\textquotedblright{}.

\begin{quote}
\te{ఆ} \te{అబ్బాయికి} \te{కొంచెం} \te{పొగరు} \te{ఎక్కువ}.
\end{quote}

\begin{quote}
\emph{ā abbāyiki koñceṁ pogaru ekkuva.}
\end{quote}

\begin{quote}
That boy is a bit too cheeky.
\end{quote}

Colloquial and blunt: \te{పొగరుగా} \te{మాట్లాడు}, to talk cheekily.

\begin{cousinweb}
\textbf{\te{పొగరు}} (\emph{pogaru}) is native Dravidian: DEDR 4234 lists Telugu \emph{pogaru} with Tamil \emph{pukalvu}, \textquotedblleft{}pride, arrogance\textquotedblright{}, and Kannada \emph{puggu}.
\end{cousinweb}

\begin{grammarlens}[title={What you've built}]
One more word for describing what you lived and felt.
\end{grammarlens}

\subsection*{Your turn}
\begin{itemize}
\raggedright
  \item \emph{Say it:} \emph{pogaru}
  \item \emph{Say it:} \emph{pogaru}, once more
  \item \emph{Say it:} \emph{ā abbāyiki koñceṁ pogaru ekkuva.} — That boy is a bit too cheeky.
  \item \emph{Recall:} say the Telugu for humility, then the Telugu for arrogance, then say \emph{pogaru} again
\end{itemize}

\begin{tcolorbox}[breakable,colback=teal!4,colframe=teal!35!black,title={Before you move on}]
What does \te{పొగరు} mean? (\textquotedblleft{}Conceit; cheek\textquotedblright{}.) Say it once more.
\end{tcolorbox}
\section[\texorpdfstring{\te{మొండితనం} (moṇḍitanaṁ)}{మొండితనం (moṇḍitanaṁ)}]{\te{మొండితనం} (moṇḍitanaṁ) — stubbornness}

\label{lesson:TE-C508-monditanam}



\begin{quote}
Before the new one: say the Telugu for arrogance, then the Telugu for conceit.
\end{quote}

\subsection*{You'll want to know: \te{మొండితనం}}
\textbf{\te{మొండితనం}} — \emph{moṇḍitanaṁ} — \textquotedblleft{}stubbornness\textquotedblright{}.

\begin{quote}
\te{పాప} \te{మొండితనంతో} \te{అన్నం} \te{తినలేదు}.
\end{quote}

\begin{quote}
\emph{pāpa moṇḍitanantō annaṁ tinalēdu.}
\end{quote}

\begin{quote}
Out of stubbornness the little girl did not eat.
\end{quote}

\te{మొండి}, stubborn, + -\te{తనం}. \te{మొండిగా}, stubbornly.

\begin{cousinweb}
\textbf{\te{మొండితనం}} (\emph{moṇḍitanaṁ}) is native Dravidian: DEDR 4971 lists Telugu \emph{moṇḍi}, \textquotedblleft{}stubborn, obstinate, perverse\textquotedblright{}, beside Tamil \emph{muraṭu}, \textquotedblleft{}roughness, obstinacy\textquotedblright{}; -\te{తనం} makes the abstract noun.
\end{cousinweb}

\begin{grammarlens}[title={What you've built}]
One more word for describing what you lived and felt.
\end{grammarlens}

\subsection*{Your turn}
\begin{itemize}
\raggedright
  \item \emph{Say it:} \emph{moṇḍitanaṁ}
  \item \emph{Say it:} \emph{moṇḍitanaṁ}, once more
  \item \emph{Say it:} \emph{pāpa moṇḍitanantō annaṁ tinalēdu.} — Out of stubbornness the little girl did not eat.
  \item \emph{Recall:} say the Telugu for arrogance, then the Telugu for conceit, then say \emph{moṇḍitanaṁ} again
\end{itemize}

\begin{tcolorbox}[breakable,colback=teal!4,colframe=teal!35!black,title={Before you move on}]
What does \te{మొండితనం} mean? (\textquotedblleft{}Stubbornness\textquotedblright{}.) Say it once more.
\end{tcolorbox}
\section[\texorpdfstring{\te{పట్టుదల} (paṭṭudala)}{పట్టుదల (paṭṭudala)}]{\te{పట్టుదల} (paṭṭudala) — determination; persistence}

\label{lesson:TE-C508-pattudala}



\begin{quote}
Before the new one: say the Telugu for conceit, then the Telugu for stubbornness.
\end{quote}

\subsection*{You'll want to know: \te{పట్టుదల}}
\textbf{\te{పట్టుదల}} — \emph{paṭṭudala} — \textquotedblleft{}determination; persistence\textquotedblright{}.

\begin{quote}
\te{పట్టుదలతో} \te{చదివి} \te{ఆమె} \te{డాక్టరు} \te{అయింది}.
\end{quote}

\begin{quote}
\emph{paṭṭudalatō cadivi āme ḍākṭaru ayindi.}
\end{quote}

\begin{quote}
By studying with determination, she became a doctor.
\end{quote}

\te{పట్టుదలతో}, with grit; \te{పట్టు}, to hold.

\begin{cousinweb}
\textbf{\te{పట్టుదల}} (\emph{paṭṭudala}) is native Dravidian: DEDR 4034 lists Telugu \emph{paṭṭu}, \textquotedblleft{}to hold, to seize\textquotedblright{}, also as a noun, \textquotedblleft{}grip, perseverance\textquotedblright{}, and \emph{paṭṭudala} itself.
\end{cousinweb}

\begin{grammarlens}[title={What you've built}]
One more word for describing what you lived and felt.
\end{grammarlens}

\subsection*{Your turn}
\begin{itemize}
\raggedright
  \item \emph{Say it:} \emph{paṭṭudala}
  \item \emph{Say it:} \emph{paṭṭudala}, once more
  \item \emph{Say it:} \emph{paṭṭudalatō cadivi āme ḍākṭaru ayindi.} — By studying with determination, she became a doctor.
  \item \emph{Recall:} say the Telugu for conceit, then the Telugu for stubbornness, then say \emph{paṭṭudala} again
\end{itemize}

\begin{tcolorbox}[breakable,colback=teal!4,colframe=teal!35!black,title={Before you move on}]
What does \te{పట్టుదల} mean? (\textquotedblleft{}Determination; persistence\textquotedblright{}.) Say it once more.
\end{tcolorbox}
\section[\texorpdfstring{\te{సాహసం} (sāhasaṁ)}{సాహసం (sāhasaṁ)}]{\te{సాహసం} (sāhasaṁ) — daring; boldness}

\label{lesson:TE-C508-sahasam}



\begin{quote}
Before the new one: say the Telugu for stubbornness, then the Telugu for determination.
\end{quote}

\subsection*{You'll want to know: \te{సాహసం}}
\textbf{\te{సాహసం}} — \emph{sāhasaṁ} — \textquotedblleft{}daring; boldness\textquotedblright{}.

\begin{quote}
\te{రాత్రి} \te{ఒంటరిగా} \te{అడవిలో} \te{నడవడం} \te{సాహసమే}.
\end{quote}

\begin{quote}
\emph{rātri oṇṭarigā aḍavilō naḍavaḍaṁ sāhasamē.}
\end{quote}

\begin{quote}
Walking alone in the forest at night is daring indeed.
\end{quote}

\te{సాహసంగా}, boldly; \te{సాహసి}, a daring person.

\begin{cousinweb}
\textbf{\te{సాహసం}} (\emph{sāhasaṁ}) is borrowed from Sanskrit \emph{sāhasa}, \textquotedblleft{}boldness; aggression\textquotedblright{} (CDIAL 13378).
\end{cousinweb}

\begin{grammarlens}[title={What you've built}]
One more word for describing what you lived and felt.
\end{grammarlens}

\subsection*{Your turn}
\begin{itemize}
\raggedright
  \item \emph{Say it:} \emph{sāhasaṁ}
  \item \emph{Say it:} \emph{sāhasaṁ}, once more
  \item \emph{Say it:} \emph{rātri oṇṭarigā aḍavilō naḍavaḍaṁ sāhasamē.} — Walking alone in the forest at night is daring indeed.
  \item \emph{Recall:} say the Telugu for stubbornness, then the Telugu for determination, then say \emph{sāhasaṁ} again
\end{itemize}

\begin{tcolorbox}[breakable,colback=teal!4,colframe=teal!35!black,title={Before you move on}]
What does \te{సాహసం} mean? (\textquotedblleft{}Daring; boldness\textquotedblright{}.) Say it once more.
\end{tcolorbox}
\section[\texorpdfstring{\te{పిరికి} (piriki)}{పిరికి (piriki)}]{\te{పిరికి} (piriki) — timid; cowardly}

\label{lesson:TE-C508-piriki}



\begin{quote}
Before the new one: say the Telugu for determination, then the Telugu for daring.
\end{quote}

\subsection*{You'll want to know: \te{పిరికి}}
\textbf{\te{పిరికి}} — \emph{piriki} — \textquotedblleft{}timid; cowardly\textquotedblright{}.

\begin{quote}
\te{మా} \te{కుక్క} \te{చాలా} \te{పిరికిది}.
\end{quote}

\begin{quote}
\emph{mā kukka cālā pirikidi.}
\end{quote}

\begin{quote}
Our dog is very timid.
\end{quote}

\te{పిరికివాడు}, a coward.

\begin{cousinweb}
\textbf{\te{పిరికి}} (\emph{piriki}) is native Dravidian: DEDR 4200 lists Telugu \emph{piṟiki}, \textquotedblleft{}timid, cowardly\textquotedblright{}, with Kannada \emph{piṟiki}, \textquotedblleft{}coward\textquotedblright{}.
\end{cousinweb}

\begin{grammarlens}[title={What you've built}]
One more word for describing what you lived and felt.
\end{grammarlens}

\subsection*{Your turn}
\begin{itemize}
\raggedright
  \item \emph{Say it:} \emph{piriki}
  \item \emph{Say it:} \emph{piriki}, once more
  \item \emph{Say it:} \emph{mā kukka cālā pirikidi.} — Our dog is very timid.
  \item \emph{Recall:} say the Telugu for determination, then the Telugu for daring, then say \emph{piriki} again
\end{itemize}

\begin{tcolorbox}[breakable,colback=teal!4,colframe=teal!35!black,title={Before you move on}]
What does \te{పిరికి} mean? (\textquotedblleft{}Timid; cowardly\textquotedblright{}.) Say it once more.
\end{tcolorbox}
